% Options for packages loaded elsewhere
\PassOptionsToPackage{unicode}{hyperref}
\PassOptionsToPackage{hyphens}{url}
\PassOptionsToPackage{dvipsnames,svgnames,x11names}{xcolor}
%
\documentclass[
  11pt,
  a4paper,
]{article}
\usepackage{amsmath,amssymb}
\usepackage{iftex}
\ifPDFTeX
  \usepackage[T1]{fontenc}
  \usepackage[utf8]{inputenc}
  \usepackage{textcomp} % provide euro and other symbols
\else % if luatex or xetex
  \usepackage{unicode-math} % this also loads fontspec
  \defaultfontfeatures{Scale=MatchLowercase}
  \defaultfontfeatures[\rmfamily]{Ligatures=TeX,Scale=1}
\fi
\usepackage{lmodern}
\ifPDFTeX\else
  % xetex/luatex font selection
\fi
% Use upquote if available, for straight quotes in verbatim environments
\IfFileExists{upquote.sty}{\usepackage{upquote}}{}
\IfFileExists{microtype.sty}{% use microtype if available
  \usepackage[]{microtype}
  \UseMicrotypeSet[protrusion]{basicmath} % disable protrusion for tt fonts
}{}
\makeatletter
\@ifundefined{KOMAClassName}{% if non-KOMA class
  \IfFileExists{parskip.sty}{%
    \usepackage{parskip}
  }{% else
    \setlength{\parindent}{0pt}
    \setlength{\parskip}{6pt plus 2pt minus 1pt}}
}{% if KOMA class
  \KOMAoptions{parskip=half}}
\makeatother
\usepackage{xcolor}
\usepackage[margin=2.25cm]{geometry}
\usepackage{color}
\usepackage{fancyvrb}
\newcommand{\VerbBar}{|}
\newcommand{\VERB}{\Verb[commandchars=\\\{\}]}
\DefineVerbatimEnvironment{Highlighting}{Verbatim}{commandchars=\\\{\}}
% Add ',fontsize=\small' for more characters per line
\newenvironment{Shaded}{}{}
\newcommand{\AlertTok}[1]{\textcolor[rgb]{1.00,0.00,0.00}{\textbf{#1}}}
\newcommand{\AnnotationTok}[1]{\textcolor[rgb]{0.38,0.63,0.69}{\textbf{\textit{#1}}}}
\newcommand{\AttributeTok}[1]{\textcolor[rgb]{0.49,0.56,0.16}{#1}}
\newcommand{\BaseNTok}[1]{\textcolor[rgb]{0.25,0.63,0.44}{#1}}
\newcommand{\BuiltInTok}[1]{\textcolor[rgb]{0.00,0.50,0.00}{#1}}
\newcommand{\CharTok}[1]{\textcolor[rgb]{0.25,0.44,0.63}{#1}}
\newcommand{\CommentTok}[1]{\textcolor[rgb]{0.38,0.63,0.69}{\textit{#1}}}
\newcommand{\CommentVarTok}[1]{\textcolor[rgb]{0.38,0.63,0.69}{\textbf{\textit{#1}}}}
\newcommand{\ConstantTok}[1]{\textcolor[rgb]{0.53,0.00,0.00}{#1}}
\newcommand{\ControlFlowTok}[1]{\textcolor[rgb]{0.00,0.44,0.13}{\textbf{#1}}}
\newcommand{\DataTypeTok}[1]{\textcolor[rgb]{0.56,0.13,0.00}{#1}}
\newcommand{\DecValTok}[1]{\textcolor[rgb]{0.25,0.63,0.44}{#1}}
\newcommand{\DocumentationTok}[1]{\textcolor[rgb]{0.73,0.13,0.13}{\textit{#1}}}
\newcommand{\ErrorTok}[1]{\textcolor[rgb]{1.00,0.00,0.00}{\textbf{#1}}}
\newcommand{\ExtensionTok}[1]{#1}
\newcommand{\FloatTok}[1]{\textcolor[rgb]{0.25,0.63,0.44}{#1}}
\newcommand{\FunctionTok}[1]{\textcolor[rgb]{0.02,0.16,0.49}{#1}}
\newcommand{\ImportTok}[1]{\textcolor[rgb]{0.00,0.50,0.00}{\textbf{#1}}}
\newcommand{\InformationTok}[1]{\textcolor[rgb]{0.38,0.63,0.69}{\textbf{\textit{#1}}}}
\newcommand{\KeywordTok}[1]{\textcolor[rgb]{0.00,0.44,0.13}{\textbf{#1}}}
\newcommand{\NormalTok}[1]{#1}
\newcommand{\OperatorTok}[1]{\textcolor[rgb]{0.40,0.40,0.40}{#1}}
\newcommand{\OtherTok}[1]{\textcolor[rgb]{0.00,0.44,0.13}{#1}}
\newcommand{\PreprocessorTok}[1]{\textcolor[rgb]{0.74,0.48,0.00}{#1}}
\newcommand{\RegionMarkerTok}[1]{#1}
\newcommand{\SpecialCharTok}[1]{\textcolor[rgb]{0.25,0.44,0.63}{#1}}
\newcommand{\SpecialStringTok}[1]{\textcolor[rgb]{0.73,0.40,0.53}{#1}}
\newcommand{\StringTok}[1]{\textcolor[rgb]{0.25,0.44,0.63}{#1}}
\newcommand{\VariableTok}[1]{\textcolor[rgb]{0.10,0.09,0.49}{#1}}
\newcommand{\VerbatimStringTok}[1]{\textcolor[rgb]{0.25,0.44,0.63}{#1}}
\newcommand{\WarningTok}[1]{\textcolor[rgb]{0.38,0.63,0.69}{\textbf{\textit{#1}}}}
\setlength{\emergencystretch}{3em} % prevent overfull lines
\providecommand{\tightlist}{%
  \setlength{\itemsep}{0pt}\setlength{\parskip}{0pt}}
\setcounter{secnumdepth}{-\maxdimen} % remove section numbering
\ifLuaTeX
  \usepackage{selnolig}  % disable illegal ligatures
\fi
\usepackage{bookmark}
\IfFileExists{xurl.sty}{\usepackage{xurl}}{} % add URL line breaks if available
\urlstyle{same}
\hypersetup{
  colorlinks=true,
  linkcolor={blue},
  filecolor={Maroon},
  citecolor={Blue},
  urlcolor={blue},
  pdfcreator={LaTeX via pandoc}}

\author{}
\date{}

\usepackage{microtype}
\usepackage{fancyhdr}
\pagestyle{fancy}
\fancyhf{}
\lhead{Documentation-Driven Threat Analysis}
\rhead{BA4 closure}
\cfoot{\thepage}
\setlength{\headheight}{14pt}
\setlength{\emergencystretch}{4em}
\hbadness=10000
\hfuzz=20pt
\sloppy
\begin{document}

\section{DDTA research work plan after documentation-layer
closure}\label{ddta-research-work-plan-after-documentation-layer-closure}

\textbf{WORK PLAN - REVISION 19}

\textbf{Status:} active plan after BA4 integrated closure.

\textbf{Supersedes:} Revision 18 only for forward execution state;
R1-R18 remain historical research records.

\subsection{1. Fixed prior state}\label{fixed-prior-state}

\begin{itemize}
\tightlist
\item
  Chapters 2-4: CLOSED / FINAL for current thesis scope.
\item
  Documentation layer: CLOSED.
\item
  BA0 responsibility/non-goals: CLOSED.
\item
  BA1 minimal BAE identity ontology: CLOSED.
\item
  BA2 relation/action vocabulary: CLOSED BY BA2-T4.
\item
  BA3 provenance/derivation/identity/lifecycle/change-revalidation:
  CLOSED BY BA3-T4.
\item
  BA4 projection boundary/traceability/coverage/interpretation: CLOSED
  BY BA4-T3.
\item
  \texttt{BAReferent} and \texttt{BAProposition}: ACCEPTED first-class
  semantic identity families.
\item
  ThreatForge remains a case study/tooling instrument, not DDTA semantic
  authority.
\end{itemize}

\subsection{2. BA4 closure result}\label{ba4-closure-result}

BA4 closes the smallest current-scope projection contract for:

\begin{itemize}
\tightlist
\item
  immutable revisioned projection definitions;
\item
  explicit eligible BA scope;
\item
  \texttt{EXHAUSTIVE\_FOR\_DECLARED\_SCOPE\ \textbar{}\ SELECTIVE}
  coverage semantics;
\item
  qualification policy for current/stale/diagnostic/pending BA inputs;
\item
  meaning-bearing projection-item trace to BA;
\item
  role-bound trace where contribution roles differ;
\item
  semantic-preserving shared rendering;
\item
  explicitly downstream method-owned interpretation;
\item
  projection-local interpretation-rule accountability;
\item
  BA trace plus BA3 continuity as the cross-projection common anchor;
  and
\item
  baseline-scoped projection rebuild without a second project lifecycle.
\end{itemize}

BA4-T3 removes independent \texttt{omissionSemantics} as redundant
because omission meaning is fixed by the coverage mode.

No universal projection ontology, universal projection DSL, new BAE
family, new BA2 operator or BA1-BA3 reopen is forced.

The active closure artifact is:

\begin{itemize}
\tightlist
\item
  \path{methodology/BA4_PROJECTION_BOUNDARY_TRACEABILITY_INTERPRETATION_COVERAGE_CONTRACT_R1.md}
\end{itemize}

BA4-T1/T2 artifacts remain immutable research derivation history.

\subsection{3. Base Analysis status after
BA4}\label{base-analysis-status-after-ba4}

\begin{Shaded}
\begin{Highlighting}[]
\NormalTok{BA0   CLOSED}
\NormalTok{BA1   CLOSED}
\NormalTok{BA2   CLOSED}
\NormalTok{BA3   CLOSED}
\NormalTok{BA4   CLOSED}
\NormalTok{BA5   NOT STARTED / NEXT}
\NormalTok{BA6   NOT STARTED}
\end{Highlighting}
\end{Shaded}

The Base Analysis milestone as a whole is not yet closed. BA5 must
establish lexical/display/assistance boundaries and BA6 must perform the
integrated corpus plus holdout closure regression.

\subsection{4. BA5 objective - lexical vocabulary and optional
assistance}\label{ba5-objective---lexical-vocabulary-and-optional-assistance}

BA5 defines the smallest boundary among source wording, stable BA
semantic keys, projection-local display labels, authoring synonyms and
optional assistance.

The phase must preserve:

\begin{Shaded}
\begin{Highlighting}[]
\NormalTok{source lexical wording}
\NormalTok{    != stable BA semantic key}
\NormalTok{    != projection{-}local display label}
\NormalTok{    != authoring synonym}
\end{Highlighting}
\end{Shaded}

The purpose is not to build a universal natural-language lexicon. It is
to prevent convenient labels, synonyms or model-generated suggestions
from silently becoming semantic authority.

NLP/LLM support, if later admitted, remains candidate-producing
assistance subject to provenance/review and cannot create project or BA
truth by itself.

\subsection{5. BA5-T1 - stable semantic key, display-label and
authoring-synonym boundary
trial}\label{ba5-t1---stable-semantic-key-display-label-and-authoring-synonym-boundary-trial}

\textbf{NEXT / NOT EXECUTED BY THIS PLAN REVISION.}

BA5-T1 must perform only the first lexical-boundary trial.

It must pressure-test at least:

\begin{enumerate}
\def\labelenumi{\arabic{enumi}.}
\tightlist
\item
  one BA semantic key rendered with several human display labels without
  changing BA identity/meaning;
\item
  one source phrase that should not become the stable BA key merely
  because it is common in the documentation;
\item
  projection-local method labels that must remain distinct from BA
  semantic keys;
\item
  authoring synonyms that can help humans express existing semantics
  without creating new semantics;
\item
  ambiguity where one lexical phrase can map to more than one existing
  BA meaning and therefore cannot be silently normalized;
\item
  whether lexical metadata requires any new first-class BAE family or
  BA2 operator; and
\item
  whether a material counterexample forces the smallest BA1-BA4 reopen.
\end{enumerate}

\subsubsection{BA5-T1 guardrails}\label{ba5-t1-guardrails}

Do not:

\begin{itemize}
\tightlist
\item
  introduce automatic arbitrary-narrative migration;
\item
  authorize NLP/LLM output as semantic truth;
\item
  define a universal domain lexicon;
\item
  convert projection-local method taxonomy into BA vocabulary;
\item
  define AnalysisRecord/Common Finding;
\item
  design ThreatForge classes/tables/UI;
\item
  reopen BA0-BA4 without a material shared-semantic counterexample;
\item
  start BA6.
\end{itemize}

\subsubsection{BA5-T1 exit condition}\label{ba5-t1-exit-condition}

Produce a falsifiable lexical responsibility lower-bound candidate and
the next smallest BA5 pressure target. Do not close BA5 from one lexical
trial.

\subsection{6. BA6 - complete Base Analysis regression and
closure}\label{ba6---complete-base-analysis-regression-and-closure}

Only after BA5 closes, regress the complete BA0-BA5 design against the
closed corpora and at least one structurally different holdout.

BA6 may reopen only the smallest earlier responsibility on a material
counterexample.

Only BA6 may close the complete Base Analysis milestone for current
thesis scope.

\subsection{7. Analysis envelope remains
downstream}\label{analysis-envelope-remains-downstream}

Only after Base Analysis closure:

\begin{itemize}
\tightlist
\item
  A1 - AnalysisRecord / execution envelope;
\item
  A2 - common reviewed Finding boundary;
\item
  A3 - change/provenance integration;
\item
  formal methodology overlays and STRIDE/STRIDE-AI evaluations.
\end{itemize}

\subsection{8. Next authorized
microstep}\label{next-authorized-microstep}

Only after the BA4-T3 closure package is reviewed, committed, pushed and
remotely verified, execute:

\begin{quote}
\textbf{BA5-T1 - stable semantic key, display-label and
authoring-synonym boundary trial.}
\end{quote}

Do not start optional assistance, BA6 or downstream analysis schemas
before BA5-T1 is completed and reviewed.

\end{document}
